\begin{abstract}
Wi-Fi sensing commonly relies on Channel State Information (CSI), but obtaining
CSI from commercial devices often requires supported hardware, modified
firmware, or specialized tools. Beamforming Feedback Information (BFI), which
is exchanged through standardized Wi-Fi beamforming procedures, provides a more
accessible alternative. However, directly applying RF-Diffusion to a
reconstructed complex beamforming matrix did not reliably preserve the
quantized angular structure, temporal beam dynamics, and activity semantics
required for BFI sensing. We present \textbf{BeamDiff}, a sensing-aware
conditional diffusion framework for generating Beamforming Feedback Angle
(BFA) sequences. BeamDiff retains the first real frame as an anchor and models
subsequent transitions as shortest-path circular deltas according to each BFA
channel's quantization range. Its objective combines diffusion noise prediction,
clean-delta reconstruction, temporal-fidelity regularization, and supervision
from a frozen activity classifier. We evaluate BeamDiff using BFI measurements
from the HAR-1 Kitchen and HAR-3 Classroom scenarios provided by the
CSI-BFI-HAR dataset. The real validation and test sets are fixed, and the
training data are augmented at a one-to-one ratio. On HAR-1, the mean BeamSense
accuracy across five model initialization seeds increases from 91.62\% to
94.63\%, a mean paired gain of 3.00 percentage points, with gains in four of
five seeds. In a paired HAR-3 evaluation, accuracy increases from 98.70\% to
99.65\%, and test errors decrease from 81 to 22. The results demonstrate that
modeling the native circular and temporal structure of BFA can produce useful
synthetic training data for BFI-based human activity sensing.
\end{abstract}
